\documentclass[11pt,a4paper]{article}

% ---------- paquetes ----------
\usepackage[spanish,es-tabla,es-noshorthands,es-nodecimaldot]{babel}
\usepackage[T1]{fontenc}
\usepackage[utf8]{inputenc}
\usepackage{lmodern}
\usepackage[margin=2.5cm]{geometry}
\usepackage{amsmath,amssymb,amsthm}
\usepackage{booktabs}
\usepackage{float}  % [H]: figuras y tablas quedan donde se escriben
\usepackage{graphicx}
\usepackage{xcolor}
\usepackage{enumitem}
\usepackage[hidelinks]{hyperref}

\graphicspath{{figs/}}
\setlist{itemsep=2pt}

% ---------- utilidades ----------
% \todo{...}: notas pendientes, en naranjo, para borrar antes de entregar
\newcommand{\todo}[1]{\textcolor{orange!80!black}{\textbf{[PENDIENTE:} #1\textbf{]}}}
% \figura{archivo}{ancho}{caption}{label}: si la imagen aún no existe, deja un recuadro
\newcommand{\figura}[4]{%
  \begin{figure}[H]
    \centering
    \IfFileExists{figs/#1}{\includegraphics[width=#2\linewidth]{#1}}%
      {\fbox{\parbox[c][4cm][c]{#2\linewidth}{\centering\small\texttt{figs/\detokenize{#1}}\\ (exportar desde el notebook)}}}
    \caption{#3}\label{#4}
  \end{figure}}

% ---------- notación ----------
\newcommand{\R}{\mathbb{R}}
\newcommand{\Sph}{\mathbb{S}^{767}}
\newcommand{\Tok}{\mathcal{V}_{\mathrm{tok}}}
\newcommand{\Img}{\mathcal{S}}
\newcommand{\Lab}{\mathcal{L}}
\newcommand{\concat}{\mathbin{\Vert}}
\newcommand{\simil}{\operatorname{sim}}
\DeclareMathOperator*{\topk}{top_{\mathit{k}}}
\DeclareMathOperator{\moda}{moda}

\title{\textbf{Mapeando YouTube fuera de la burbuja}\\[4pt]
  \large Embeddings multimodales de videos con EmbeddingGemma~2}
\author{Joaquín Bächler - Matías Morales - Franco Moya - Cristóbal Vargas \\ IIC2433 Minería de Datos — 2026-2}
\date{\today}

\begin{document}
\maketitle

\begin{abstract}
\todo{Escribir al final (5--6 líneas): problema, datos, método, resultado principal.}
Representamos $\approx 36.000$ videos de YouTube en un espacio vectorial común usando \texttt{EmbeddingGemma~2},
a partir de su texto (título, tags y descripción) y de su miniatura. Sobre el texto, el modelo
supera a dos baselines (TF-IDF y MiniLM) en todas las métricas. Sobre los videos, la imagen (thumbnail) por sí sola es casi tan informativa como el texto, y su fusión tardía mejora la
recuperación de la categoría temática (kNN de 0,798 a 0,838). Además, como imagen y texto
comparten el espacio, la miniatura permite etiquetar videos sin entrenamiento (zero-shot):
en el 60\,\% de los casos, alguna de las 5 etiquetas predichas es una etiqueta real.
\end{abstract}

% =====================================================================
\section{Introducción}
Las plataformas, como \texttt{YouTube} recomiendan contenido cercano a lo que ya consumimos, llamemos a esto \textit{burbujas de filtro}. Nos interesaría saber si ¿podemos construir un mapa del contenido de YouTube que
permita recomendar videos \emph{fuera} de esa zona? Contribuciones:
(1) pipeline de datos YT-8M + API, (2) comparación de representaciones de texto,
imagen y multimodales, (3) etiquetado zero-shot desde la miniatura.

% =====================================================================
\section{Datos}

\subsection{YouTube-8M}
Usamos la versión~2 (2018) de YouTube-8M~\cite{yt8m}, con $\sim$6,1 millones de videos.
Del split de entrenamiento descargamos una muestra aleatoria de $\sim$1\,\%: 41 archivos
(\emph{shards}), seleccionados de forma determinística por el hash MD5 de su nombre.

De ahora en adelante, se propone modelar matemáticamente este proceso, con el objetivo de darle algún sustento, orden y referencia. 

Sea $\mathcal{I} = \{1, \dots, N\}$ el conjunto de videos de la muestra, con $N = 41.393$.
Construimos la tabla $Y$, cuya fila $i$ es
\[
  y_i = (p_i,\ \ell_i,\ r_i,\ a_i),
\]
donde
\begin{itemize}
  \item $p_i \in \Sigma^4$ es un identificador anonimizado (\texttt{pid}) de 4 caracteres;
  \item $\ell_i \subseteq \Lab = \{0, \dots, 3861\}$ es el conjunto de etiquetas temáticas del
        video, con $|\ell_i| \approx 3$ en promedio;
  \item $r_i \in \R^{1024}$ es el promedio temporal de las activaciones de una red Inception
        sobre los frames del video (\texttt{mean\_rgb});
  \item $a_i \in \R^{128}$ es su análogo para el audio (\texttt{mean\_audio}).
\end{itemize}
A cada video le asignamos una \emph{categoría gruesa}
\[
  c_i = \moda\{\, \mathrm{Vertical1}(j) : j \in \ell_i \,\},
\]
con 25 valores posibles. Esta categoría es la referencia para evaluar las representaciones.
Luego nos encontraremos (y trataremos de resolver) un problema de proporcionalidad dentro de estas categorías (hay 2 que están sobre representadas en el dataset y/o abarcan-generalizan contenido).

\subsection{Metadatos desde la YouTube Data API}
Los \texttt{pid} no son identificadores de YouTube. Existe una función \emph{parcial}
$\varphi$ (una tabla de traducción publicada por Google) tal que $\varphi(p_i) = v_i$ es el
identificador real del video. Con $v_i$ consultamos la YouTube Data API~v3 y obtenemos
\[
  w_i = (t_i,\ d_i,\ g_i,\ m_i,\ u_i),
\]
con título $t_i$, descripción $d_i$, tags $g_i$, estadísticas $m_i$ (vistas, likes,
comentarios) y la URL de la miniatura $u_i$. A diferencia de $Y$, la tabla $W$ no es un
espacio vectorial: sus columnas mezclan texto, números y URLs.

No todos los videos sobreviven a estos pasos:
\[
  \mathcal{I} \;\supseteq\; \mathcal{I}_\varphi = \{ i : \varphi(p_i) \text{ está definida} \}
  \;\supseteq\; \mathcal{I}_W = \{ i : w_i \text{ existe} \}.
\]

\begin{table}[H]
  \centering
  \caption{Embudo de datos.}
  \label{tab:embudo}
  \begin{tabular}{lrr}
    \toprule
    Etapa & Videos & \% \\
    \midrule
    Muestra YT-8M ($\mathcal{I}$)                  & 41.393 & 100,0 \\
    Con ID real de YouTube ($\mathcal{I}_\varphi$) & 37.230 &  89,9 \\
    Disponibles en YouTube ($\mathcal{I}_W$)       & \textbf{34.788} & \textbf{84,0} \\
    \bottomrule
  \end{tabular}
\end{table}

En adelante trabajamos sobre $\mathcal{I}_W$, con $n = |\mathcal{I}_W| = 34.788$.

% =====================================================================
\section{Método}

\subsection{El espacio de embeddings}
\texttt{EmbeddingGemma~2}~\cite{embgemma2} es un modelo multimodal: recibe secuencias que pueden
mezclar \emph{tokens} de texto e imágenes, y las lleva a un mismo espacio. Formalmente, sea
$\Tok$ el vocabulario de tokens del modelo y $\Img$ el conjunto de imágenes RGB. El dominio del
modelo es el conjunto de secuencias finitas de tokens y/o imágenes, y su codominio es la esfera
unitaria:
\[
  f_\theta : (\Tok \sqcup \Img)^{\le 8192} \longrightarrow \Sph
  = \{\, x \in \R^{768} : \|x\|_2 = 1 \,\}.
\]
Un texto es el caso particular de una secuencia con solo tokens; una imagen, el de una
secuencia con un solo elemento de $\Img$. Como todos los vectores tienen norma~1, la
similitud coseno se reduce al producto punto:
\[
  \simil(x, x') = \frac{x^\top x'}{\|x\|\,\|x'\|} = x^\top x'.
\]

\paragraph{Texto.}
Definimos $\tau(w_i)$ como el título, seguido de hasta 15 tags y los primeros 600 caracteres
de la descripción sin URLs. El modelo recibe además un prefijo fijo que indica la tarea,
$\pi_{\mathrm{clust}} = $ \texttt{"task: clustering | query: "}. Con $\concat$ la
concatenación de strings,
\[
  x_i^{T} = f_\theta\big(\pi_{\mathrm{clust}} \concat \tau(w_i)\big), \qquad
  X^{T} = \begin{bmatrix} x_1^{T} \\ \vdots \\ x_n^{T} \end{bmatrix} \in \R^{n \times 768}.
\]
El texto se trunca a 256 tokens, lo que afecta al 2,9\,\% de los videos.

\paragraph{Imagen.}
Sea $s_i \in \{0, \dots, 255\}^{180 \times 320 \times 3}$ la miniatura del video. Sobre
$\mathcal{I}_S \subset \mathcal{I}_W$, los videos cuya miniatura se pudo descargar
($|\mathcal{I}_S| = 34.771$),
\[
  x_i^{I} = f_\theta(s_i), \qquad X^{I} \in \R^{34771 \times 768}.
\]
Todas las comparaciones de la Sección~\ref{sec:resultados} se hacen sobre $\mathcal{I}_S$, para
que cada representación se mida sobre exactamente los mismos videos.
Usamos 70 tokens visuales por imagen (el valor por defecto es 280), lo que es 5 veces más
rápido y cambia poco el vector (similitud 0,986 con la versión de 280 tokens).

\paragraph{Combinaciones de texto e imagen.}
Probamos dos formas de juntar ambas fuentes:
\[
  \underbrace{x_i^{IT} = f_\theta\big(\pi_{\mathrm{clust}} \concat s_i \concat \tau(w_i)\big)}_{\text{fusión temprana: un solo input}}
  \qquad\qquad
  \underbrace{x_i^{F} = \frac{x_i^{T} + x_i^{I}}{\big\|x_i^{T} + x_i^{I}\big\|_2}}_{\text{fusión tardía: suma de vectores}}
\]

\subsection{Baselines de texto}
Para justificar el uso de EmbeddingGemma~2, embebimos el mismo texto $\tau(w_i)$ con:
\begin{itemize}
  \item \textbf{TF-IDF + SVD}: vectores de frecuencia de palabras (vocabulario de 50.000
        términos) comprimidos a 768 dimensiones con SVD truncado (LSA).
  \item \textbf{MiniLM}~\cite{sbert,sbert-multi}: modelo Sentence-BERT multilingüe\footnote{\texttt{sentence-transformers/paraphrase-multilingual-MiniLM-L12-v2}.},
        de 384 dimensiones y 128 tokens
        como máximo.
\end{itemize}

\subsection{Zero-shot desde la miniatura}
Sea $L \in \R^{3806 \times 768}$ la matriz cuya fila $j$ es $f_\theta(\mathrm{nombre}_j)$, el
embedding del nombre de la etiqueta $j$ como texto. Como $X^{I}$ y $L$ viven en el mismo
espacio,
\[
  S = X^{I} L^{\top} \in \R^{34771 \times 3806}, \qquad
  S_{ij} = \simil(\text{imagen}_i,\ \text{etiqueta}_j).
\]
La predicción para la imagen $i$ son las $k$ etiquetas de mayor similitud, y medimos si
alguna es correcta:
\[
  \hat{\ell}_i^{(k)} = \topk_{j}\, S_{ij}, \qquad
  \text{hit@}k = \frac{1}{|\mathcal{I}_S|} \sum_{i \in \mathcal{I}_S}
  \mathbf{1}\big[\, \hat{\ell}_i^{(k)} \cap \ell_i \neq \varnothing \,\big].
\]

\subsection{Proyección y evaluación}
Para visualizar usamos UMAP~\cite{umap} $U : \Sph \to \R^3$ (métrica coseno, 15 vecinos),
opcionalmente precedido de PCA a 282 componentes, que retienen el 90\,\% de la varianza.
Cada representación $X$ se evalúa con:
\begin{itemize}
  \item \textbf{kNN accuracy}: qué tan bien los 15 vecinos más cercanos de $x_i$ predicen
        $c_i$ (validación cruzada de 5 folds), en el espacio original y en $U(X)$;
  \item \textbf{NMI}: información mutua normalizada entre los clusters de K-Means ($k = 25$)
        sobre $U(X)$ y las categorías $c$;
  \item \textbf{Trustworthiness}~\cite{trust}: qué tan fieles son los vecindarios de $U(X)$
        a los de $X$.
\end{itemize}
Como referencia mínima usamos el clasificador de mayoría (siempre predecir la categoría más
frecuente).

% =====================================================================
\section{Resultados}
\label{sec:resultados}

\subsection{Texto: EmbeddingGemma 2 frente a los baselines}

\begin{table}[H]
  \centering
  \caption{Representaciones de texto sobre los 34.771 videos de $\mathcal{I}_S$.}
  \label{tab:texto}
  \small
  \begin{tabular}{lrrrrr}
    \toprule
    Modelo & Dims & kNN original & kNN UMAP & NMI (UMAP) & Trustworthiness \\
    \midrule
    TF-IDF + SVD     & 768 & 0,609 & 0,560 & 0,134 & 0,863 \\
    MiniLM           & 384 & 0,698 & 0,675 & 0,297 & 0,888 \\
    EmbeddingGemma 2 & 768 & \textbf{0,798} & \textbf{0,771} & \textbf{0,378} & \textbf{0,956} \\
    \bottomrule
  \end{tabular}
\end{table}

\todo{Comentar: el orden se mantiene en todas las métricas; +19 puntos de kNN sobre TF-IDF
y +10 sobre MiniLM; NMI casi triplica a TF-IDF.}

\figura{metricas_modelos.png}{0.95}{Comparación de modelos de embedding de texto.}{fig:metricas-modelos}

\begin{table}[H]
  \centering
  \caption{Proyecciones de EmbeddingGemma 2 (texto).}
  \label{tab:proyecciones}
  \begin{tabular}{lrrrr}
    \toprule
    Proyección & Trustworthiness & Silhouette & kNN & NMI \\
    \midrule
    PCA 3D              & 0,828 & $-0{,}133$ & 0,643 & 0,325 \\
    UMAP 3D             & 0,956 & $-0{,}099$ & 0,771 & \textbf{0,378} \\
    PCA 282 + UMAP 3D   & \textbf{0,958} & $-0{,}101$ & 0,773 & 0,364 \\
    Original (768 dims) & --    & $-0{,}003$ & \textbf{0,798} & 0,351 \\
    Mayoría             & --    & --         & 0,347 & -- \\
    \bottomrule
  \end{tabular}
\end{table}

\todo{PCA en 3D retiene solo el 11,2\,\% de la varianza. Explicar el silhouette negativo:
coherencia local (kNN) sin separación global de categorías gruesas.}

\figura{metricas_text.png}{0.95}{Calidad de las proyecciones de EmbeddingGemma 2.}{fig:metricas-texto}

\subsection{Texto, imagen y fusión}

\begin{table}[H]
  \centering
  \caption{Comparación de las seis representaciones sobre los 34.771 videos de $\mathcal{I}_S$.}
  \label{tab:multimodal}
  \begin{tabular}{lrrr}
    \toprule
    Representación & kNN original & kNN UMAP & NMI (UMAP) \\
    \midrule
    Mayoría                       & 0,347 & -- & -- \\
    TF-IDF + SVD                  & 0,609 & 0,560 & 0,134 \\
    MiniLM                        & 0,698 & 0,675 & 0,297 \\
    Texto, $X^{T}$                & 0,798 & 0,771 & 0,378 \\
    Imagen, $X^{I}$               & 0,784 & 0,769 & 0,391 \\
    Fusión temprana, $X^{IT}$     & 0,819 & 0,791 & 0,364 \\
    Fusión tardía, $X^{F}$        & \textbf{0,838} & \textbf{0,823} & \textbf{0,452} \\
    \bottomrule
  \end{tabular}
\end{table}

\todo{Comentar: la miniatura sola casi iguala al texto; la fusión tardía es la mejor en las tres
métricas y no cuesta cómputo extra; la temprana mejora al texto, pero bastante menos que la tardía
(hipótesis: el texto domina el input).}

\figura{metricas_multimodal.png}{0.95}{Texto frente a imagen y fusión.}{fig:metricas-multimodal}
\figura{vecinos.png}{0.95}{Vecinos más cercanos de un video en los espacios de texto, imagen y fusión.}{fig:vecinos}

\subsection{Zero-shot}

\begin{table}[H]
  \centering
  \caption{Etiquetado zero-shot desde la miniatura (3.806 etiquetas candidatas).}
  \label{tab:zeroshot}
  \begin{tabular}{lrrr}
    \toprule
    & hit@1 & hit@5 & hit@10 \\
    \midrule
    Imagen, $X^{I}$ & 32,2\,\% & 60,0\,\% & 70,2\,\% \\
    \bottomrule
  \end{tabular}
\end{table}

\figura{zeroshot_hitk.png}{0.7}{Fracción de videos con al menos una etiqueta real entre las $k$ predichas.}{fig:hitk}
\figura{zeroshot_galeria.png}{0.95}{Ejemplos de etiquetado zero-shot (\checkmark{} = etiqueta real).}{fig:zeroshot}

% =====================================================================
\section{Discusión}
\todo{Desarrollar cada punto:}
\begin{itemize}
  \item \textbf{Brecha de modalidad}~\cite{modalitygap}: la similitud media imagen--imagen es
        0,67 y la imagen--texto, 0,55; en un UMAP conjunto forman dos islas separadas, que se
        juntan al centrar cada modalidad.
  \item \textbf{Sesgo por idioma}: TF-IDF agrupa por idioma más que por tema.
  \item \textbf{Fusión tardía frente a temprana}.
\end{itemize}

\section{Limitaciones}
\begin{itemize}
  \item Las etiquetas de YT-8M fueron generadas por el sistema de anotación de YouTube a partir
        de ``señales humanas, incluyendo metadatos y clics de búsqueda''~\cite{yt8m}. Como en parte
        provienen del texto del video, la evaluación podría favorecer a los embeddings de texto
        frente a los de imagen.
  \item La categoría $c_i$ es gruesa (25 valores) y es una sola por video, aunque el video tenga
        varias etiquetas: un video de un juego de carreras puede ser \emph{Games} o
        \emph{Autos \& Vehicles}. Las métricas basadas en $c_i$ son, por lo tanto, una cota
        inferior de la calidad de las representaciones.
  \item Truncación de texto (256 tokens en Gemma~2, 128 en MiniLM).
  \item El 16\,\% de los videos ya no está disponible, lo que puede sesgar la muestra.
\end{itemize}

\section{Trabajo futuro}
\begin{itemize}
  \item HDBSCAN sobre UMAP y sensibilidad a \texttt{n\_neighbors}.
  \item Usar la fusión tardía, la mejor representación medida, en el mapa de la aplicación web.
  \item Recomendador \emph{fuera de la burbuja}, evaluado con novedad, diversidad y cobertura.
  \item Análisis causal: efecto del tema o de la miniatura sobre las vistas.
\end{itemize}

% =====================================================================
\begin{thebibliography}{9}
\bibitem{yt8m}
  S.~Abu-El-Haija \emph{et al.}
  YouTube-8M: A Large-Scale Video Classification Benchmark.
  \emph{arXiv:1609.08675}, 2016.

\bibitem{embgemma2}
  Google DeepMind.
  EmbeddingGemma 2 (model card).
  \url{https://huggingface.co/google/embeddinggemma-2}.

\bibitem{sbert}
  N.~Reimers e I.~Gurevych.
  Sentence-BERT: Sentence Embeddings using Siamese BERT-Networks.
  \emph{EMNLP}, 2019.

\bibitem{sbert-multi}
  N.~Reimers e I.~Gurevych.
  Making Monolingual Sentence Embeddings Multilingual using Knowledge Distillation.
  \emph{EMNLP}, 2020.

\bibitem{umap}
  L.~McInnes, J.~Healy y J.~Melville.
  UMAP: Uniform Manifold Approximation and Projection for Dimension Reduction.
  \emph{arXiv:1802.03426}, 2018.

\bibitem{trust}
  J.~Venna y S.~Kaski.
  Neighborhood Preservation in Nonlinear Projection Methods: An Experimental Study.
  \emph{ICANN}, 2001.

\bibitem{modalitygap}
  W.~Liang \emph{et al.}
  Mind the Gap: Understanding the Modality Gap in Multi-modal Contrastive Representation Learning.
  \emph{NeurIPS}, 2022.
\end{thebibliography}

\end{document}
